% !PRINTED 250
\weberpage{250}
%<<CONT
erstes Glied ein Quadrat und deren letztes ein Nichtquadrat ist, ein quadratisches Glied auf, auf welches ein Nichtquadrat folgt.

Ist dann also $a = \xi^2$ ein Quadrat und $\xi^2 + 1$ ein Nichtquadrat, so ist, wenn $\beta$ ein Nichtquadrat ist, $\beta : \xi^2 + 1 = \eta^2$ ein Quadrat, und daraus folgt:
\[ \beta = \xi^2 \eta^2 + \eta^2 \]
w. z. b. w.

\begin{center}
\S.\,65.

\gesperrt{Congruenzgruppen im K\"orper $\mathfrak{C}$.}
\end{center}

In jedem Congruenzk\"orper $\mathfrak{C}$ kann man eine endliche Gruppe von linearen Substitutionen bilden, wenn man in
\weq{1}{A = \begin{pmatrix} \alpha, & \beta \\ \gamma, & \delta \end{pmatrix}}
die Elemente $\alpha, \beta, \gamma, \delta$ alle Zahlen von $\mathfrak{C}$ durchlaufen l\"asst, und wenn man je zwei dieser Substitutionen nach der Vorschrift
\weq{2}{\begin{pmatrix} \alpha, & \beta \\ \gamma, & \delta \end{pmatrix} \begin{pmatrix} \alpha', & \beta' \\ \gamma', & \delta' \end{pmatrix} = \begin{pmatrix} \alpha\alpha' + \beta\gamma', & \alpha\beta' + \beta\delta' \\ \gamma\alpha' + \delta\gamma', & \gamma\beta' + \delta\delta' \end{pmatrix}}
componirt. Der Grad dieser Gruppe ist $p^{4n}$, wenn $p$ der Modul und $n$ der Grad von $\mathfrak{C}$ ist.

Darin ist aber eine Gruppe niedrigeren Grades enthalten. Nennen wir $\alpha\delta - \beta\gamma$ die Determinante der Substitution (1), so folgt aus (2), dass die Determinante der aus $A$ und $A'$ componirten Substitution $AA'$ gleich dem Producte der Determinanten von $A$ und $A'$ ist. Wenn wir daher den Zahlen $\alpha, \beta, \gamma, \delta$ die Bedingung auferlegen, dass
\weq{3}{\alpha\delta - \beta\gamma = 1}
sein soll, so bilden auch diese Elemente $A$ eine Gruppe. Der Grad dieser Gruppe ist gleich der Anzahl der L\"osungen von (3). Um diese Anzahl zu finden, bemerken wir zun\"achst, dass wir f\"ur $\alpha, \beta$ irgend zwei Zahlen aus $\mathfrak{C}$ setzen k\"onnen, die nicht beide gleich Null sind. Denn ist etwa $\alpha$ von Null verschieden, so kann man $\gamma = 0, \delta = 1 : \alpha$ setzen und hat so die Bedingung (3) erf\"ullt. Die Anzahl der brauchbaren Zahlenpaare $\alpha, \beta$ ist also $p^{2n} - 1$. Hat man aber zu einem Zahlenpaare $\alpha, \beta$ eine L\"osung $\gamma_0, \delta_0$ von (3) gefunden, so m\"ussen alle \"ubrigen der Bedingung
\[ \alpha (\delta - \delta_0) = \beta (\gamma - \gamma_0) \]
%CONT>>
